Latour, \textit{Jamais fomos modernos}.
Na p. 64, diz o autor que \enquote{o pós-modernismo rejeita qualquer trabalho empírico como sendo ilusório e enganador. Racionalistas decepcionados, seus adeptos sentem claramente que o modernismo terminou, mas continuam a aceitar sua forma de dividir o tempo e não podem, portanto, recortar as épocas senão através de revoluções que se sucederiam umas às outras (...) \textit{No future}, é o seu slogan que acrescenta-se ao dos modernos, \textit{No past}.} A crítica despatologizante contra a memificação das diferenças inverte o sinal da violência diagnóstica apontada quando nega a possibilidade de diferenciação positiva dos indivíduos diagnosticados tendo em vista que os prejuízos enfrentados por tais indivíduos decorreria não de diferenças encontradas em seus corpos, mas de vivermos segundo arranjamentos sócio-econômicos patológicos eles próprios.